\section{How a high zero tail can hide the continuous defect (NS-77)}
\label{sec:ns77-continuous-tail}

\paragraph{Scope.}
We give a complete upper bound on the continuous defect in NS-76 using
positive sums over hypothetical off-critical zeros. This quantifies a
limitation of finite geometric evidence. It is not an upper bound on
$E_N$ or its integer limit, not a new verified zero height, and not a
proof that the defect is exactly zero.

Suppose every off-critical zero in the upper half-plane has height
$\gamma>H\geq1$. Index the conjugate pairs by $\rho=\beta+i\gamma$,
with $1/2<\beta<1$, once per pair and with multiplicity, and put
\[
 S_2=\sum_\rho\frac{2\beta-1}{\gamma^2},\qquad
 S_4=\sum_\rho\frac{2\beta-1}{\gamma^4}.
\]
These sums converge by the usual zero-counting bound. Empty sums are zero.

\begin{proposition}[A complete positive zero-tail budget]
\label{prop:ns77-tail-budget}
For the continuous $q=2$ distance,
\begin{equation}
 0\leq D_{\rm cont}^2\leq
       \min\left\{2,\frac83S_2^3+5S_4\right\}.
 \label{eq:ns77-tail-budget}
\end{equation}
If the total upper-half-plane zero count satisfies
$N(t)\leq A t\log t$ for all $t\geq H\geq100$, then
\begin{equation}
 D_{\rm cont}^2\leq
 \frac{\frac{64}{3}A^3(\log H+1)^3
                  +20A(\log H/3+1/9)}{H^3}.
 \label{eq:ns77-height-bound}
\end{equation}
No tail prefix is substituted for the complete zero sums.
\end{proposition}
\begin{proof}
For one pair write $a=\beta$, $c=1-\beta$, $\delta=a-c=2\beta-1$,
so $a+c=1$. Its real-normalized inner factor at $1$ is
$B_\rho(1)=(\gamma^2+c^2)/(\gamma^2+a^2)$.
Set $j_\rho=-\log B_\rho(1)$ and
$\ell_\rho=(\log B_\rho)'(1)$. Exactly,
\[
 j_\rho=\int_{c^2}^{a^2}\frac{dt}{\gamma^2+t},\qquad
 \ell_\rho=-\frac{2\delta(\gamma^2-ac)}
                        {(\gamma^2+a^2)(\gamma^2+c^2)}.
\]
Thus $0\leq j_\rho\leq\delta/\gamma^2$ and
\[
 \ell_\rho+\frac{2\delta}{\gamma^2}
 =\frac{2\delta\{\gamma^2(a^2+c^2+ac)+a^2c^2\}}
        {\gamma^2(\gamma^2+a^2)(\gamma^2+c^2)}
 \leq\frac{17\delta}{8\gamma^4},
\]
using $a^2+c^2+ac\leq1$, $a^2c^2\leq1/16$, and $\gamma\geq1$.
In particular $\ell_\rho+2j_\rho\leq17\delta/(8\gamma^4)$.
Absolute convergence permits summing these inequalities and the
logarithmic derivatives. With
\[
 J=-\log B(1)\in[0,S_2],\qquad \ell=B'(1)/B(1),
 \qquad \epsilon=\ell+2J,
\]
we have $\epsilon\leq17S_4/8$. There is no lower-sign assumption on
$\epsilon$.

Let $\ell_2=(\log B)''(1)$, a real number. Substituting
$b=e^{-J}$, $d=b\ell$, $e=b(\ell^2+\ell_2)$ into
\eqref{eq:ns76-continuous-distance} and discarding only the
nonnegative square that is subtracted gives
\begin{align*}
 D_{\rm cont}^2
 &\leq 2-e^{-2J}(2-2\ell+\ell^2)\\
 &=2-e^{-2J}(2+4J+4J^2)
       +(2+4J)e^{-2J}\epsilon-e^{-2J}\epsilon^2.
\end{align*}
The first term on the last line is zero at $J=0$ and has derivative
$8J^2e^{-2J}$, so it is at most $8J^3/3$.
Also $(2+4J)e^{-2J}\leq2$ for $J\geq0$. If $\epsilon\leq0$ its
linear contribution is nonpositive; otherwise it is at most
$17S_4/4<5S_4$. This proves~\eqref{eq:ns77-tail-budget}.

For the second assertion, $0<2\beta-1<1$ and partial summation give,
for $k=2,4$,
\[
 S_k\leq k\int_H^\infty\frac{N(t)}{t^{k+1}}\,dt
 \leq kA H^{1-k}
             \left(\frac{\log H}{k-1}+\frac1{(k-1)^2}\right).
\]
The endpoint at infinity vanishes and dropping the nonpositive lower
endpoint term only enlarges the bound. Inserting these complete
integrals proves~\eqref{eq:ns77-height-bound}.
\end{proof}

\paragraph{Calibration using published finite-height information.}
The bound of Hasanalizade--Shen--Wong
\cite[Corollary 1.2]{HasanalizadeShenWong2021} implies
$N(t)\leq t\log t$ for $t\geq100$. Indeed its main term is at most
$(t/6)\log t$, and its remainder is at most
$.362\log t+10$; since $\log100>4$, their sum is less than
$(t/6+3)\log t\leq t\log t$ on this range.
Platt--Trudgian's published interval verification
\cite[Theorem 1]{PlattTrudgian2021} supplies the absence of off-critical
zeros up to $H=3\cdot10^{12}$. Taking $A=1$ in the analytic bound
therefore gives the explicit consequence
\begin{equation}
 0\leq D_{\rm cont}^2<2.08\cdot10^{-32}.
 \label{eq:ns77-published-calibration}
\end{equation}
Only the final scalar evaluation is replayed here at 256/384 bits;
the cited large zero verification is an external published input,
not a new computation or an archived verification by this project.
This bound is deliberately conservative.

A small positive continuous defect remains compatible with this
inequality. Moreover the integer and continuous closures are strictly
different by Proposition~\ref{prop:ns76-strict-closure}, and no matching
upper bound on the integer limiting error has been supplied.
Thus the calibration explains why the relaxed problem can look
extremely favorable without proving the original approximation target,
a new zero-free region, G2 or RH.
